\documentclass[journal]{IEEEtai}

\usepackage[colorlinks,urlcolor=blue,linkcolor=blue,citecolor=blue]{hyperref}

\usepackage{amsmath}
\usepackage{amssymb}

\usepackage{color,array}

\usepackage{graphicx}

%% change the following later - leftover from template
%% \jvol{XX}
%% \jnum{XX}
%% \paper{1234567}
%% \pubyear{2020}
%% \publisheddate{xxxx 00, 0000}
%% \currentdate{xxxx 00, 0000}
%% \doiinfo{TQE.2020.Doi Number}


\setcounter{page}{1}
%% \setcounter{secnumdepth}{0}


\begin{document}


\title{Fine-Tuning for Age-Appropriate Pediatric Hospital Communication via LoRA}


\author{authors go here}

\maketitle

\begin{abstract}
Pediatric patients navigating hospital environments require age-appropriate, institutionally grounded information that general-purpose AI systems are not designed to provide. We present an age-targeted fine-tuning framework for a virtual hospital guide deployed within a digital recreation of Victoria Hospital (London, Ontario), specialized for children aged 5--11 and adolescents aged 12--18. Using Low-Rank Adaptation (LoRA) on 4-bit quantized Llama 3.2 Instruct models (1B and 3B) running on an Apple Mac Mini M4 (16~GB) --- a deliberate constraint modeling resource-limited clinical deployment --- we train two adapter configurations: Standard ($r=8$, 16 layers) and Fast ($r=4$, 8 layers), across both model sizes and age groups for 8 training runs total. Training data consists of approximately 475 synthetic examples per adapter generated by Claude Opus 4.6 using Victoria Hospital's public patient materials as grounding artifacts, following a LIMA-inspired curation approach \cite{zhou2023lima}. Fine-tuned adapters substantially outperform zero-shot base models on readability (72--84\% of 5--11 responses meeting a Flesch-Kincaid $\leq 7.0$ grade-level target vs.\ 12--14\% for base models) and inter-role register separation (TF-IDF + logistic regression classifier accuracy 0.89--0.94 vs.\ 0.66--0.70). A configuration-ordering crossover is observed: the Standard adapter yields better results on the 3B model while the Fast adapter performs better on the 1B model, suggesting optimal adapter capacity is model-size-dependent. Only Fast-1B meets the 1.0-second real-time VR latency target (0.93s average). These results demonstrate that compact, parameter-efficient adaptation on consumer hardware is feasible for pediatric clinical NLP.
\end{abstract}

\section{Introduction}

\IEEEPARstart{T}{his} document is a template... (intro here). Gotta get to 6 pages with the double column

\section{Related Work}
a
\subsection{Subsection}
a2
\subsection{Subsection}
a3

\section{Dataset}
The dataset used to train the models for this project was generated by Claude Opus 4.6 in extended thinking mode. (talk more about extended thinking mode here?) Claude was provided with the Children's Hospital Patient and Family Guide (cite), which contains information about the hospital, its routines, procedures, and answers questions that new patients and caretakers may have about their time at the hospital.\par

\medskip From this guide, question and answer pairs were created for five different categories. These categories include Emotional Reassurance, FAQs, Hospital Rules/Routines, What to Expect, and "Who Are These People?". (maybe go into more detail here, move the point form section talking about this in the methodology section to here?) Questions and age-appropriate answers were generated for each of the different age groups, ages 5-11 and 12-18.\par

\medskip (talk about why generated here - take talking points from methodology section? Or leave that there?)

\medskip For each role and category combination, 100 question and answer pairs were created. For example, the "Emotional Reassurance" category would have 200 questions total, 100 for the 5-11 age group, and 100 for the 12-18 age group. With five different categories, each having entries for both roles, this adds up to 1000 examples total in our dataset. \par

\medskip These 1000 examples are placed in five different .jsonl files, each corresponding to one of the five categories. Each .jsonl file includes all 200 examples for its respective category. \par

\medskip These examples are formatted as follows: 

\begin{verbatim}
{"instruction": "...", "response": "...", 
"role": "5-11|12-18", "category": "..."}
\end{verbatim}\par

\medskip The "instruction" key holds the user's question, and the "response" key holds the appropriate answer to this question. The "role" key denotes the age group of the user asking the question- this will always be "5-11" or "12-18". Lastly, the "category" key holds the category of the question asked, which would be one of the five categories mentioned prior.\par

\medskip One important thing to note in this dataset is the lack of a system prompt. This is a deliberate choice when desigining both the dataset and the project at large- as this guide will be part of the larger Virtual Hospital project, we must ensure it will fit smoothly with the rest of the project. As the team working on the VR and spatial navigation aspects of the Virtual Hospital already relies on their own system prompt for navigation, refraining from including our own system prompt avoids any potential conflicts. \par

\medskip In addition to this, when using a system prompt, we risk the model simply learning the prompt, and becoming overly reliant on it. When the system prompt is removed, this could lead to a degradation in both quality and the needed tone in the model's responses (get citation for this). By avoiding the use of a system prompt in the first place, we can ensure that the model is learning the needed tone in its weights, and not just following the prompt. (word this last sentence better?) \par

\section{Methodology}\label{sec:methodology}

\subsection{System Overview}

This work extends a virtual hospital framework designed to support safe, naturalistic communication between pediatric patients, caregivers, and clinical decision-makers. The platform instantiates a digital recreation of Victoria Hospital (London, Ontario) in which users interact with an AI-powered guide. The conversational interface routes each query to one of two age-targeted response modules: one serving younger children (ages 5--11) and one serving adolescents (ages 12--18), determined by the established user profile. Emotional-support queries are in scope only to the extent of age-appropriate reassurance and redirection; clinical assessment, therapeutic guidance, and crisis intervention are explicitly out of scope.

\subsection{Base Models and Infrastructure}

We use the Llama 3.2 Instruct model family \cite{dubey2024llama3} as our foundation, specifically the 1B and 3B parameter variants, loaded in 4-bit quantized form via the \texttt{mlx\_lm} framework \cite{apple2023mlx} on Apple Silicon. All training and evaluation runs on an Apple Mac Mini M4 (16 GB unified memory) --- a deliberate constraint modeling deployment in resource-limited clinical settings without cloud infrastructure. A lower-bound estimate for full-parameter fine-tuning of the 3B model ($\sim$30 GB in bfloat16 before activations) exceeds the 16 GB ceiling by nearly 2$\times$; the 1B model is borderline with no headroom. Parameter-efficient adaptation via LoRA is therefore the only viable fine-tuning strategy on this hardware for both model scales.

\subsection{LoRA as Declarative Programming}

We employ Low-Rank Adaptation (LoRA) \cite{hu2022lora} to specialize the base models without modifying full model weights. A pre-trained weight matrix $W_0 \in \mathbb{R}^{d \times k}$ is held frozen; a low-rank update is expressed as:

\begin{equation}
\Delta W = BA
\end{equation}

where $B \in \mathbb{R}^{d \times r}$ and $A \in \mathbb{R}^{r \times k}$, with rank $r \ll \min(d, k)$. We insert adapters into all linear projection modules within the targeted transformer layers, using \texttt{scale = 20.0} held constant across all runs.

We conceptualize LoRA adapters as a form of \textit{declarative programming} over foundation models: hospital guidelines, care pathways, and age-appropriate communication standards collectively define a behavioral specification that the adapter realizes during inference, without encoding knowledge as explicit procedural rules.

\textbf{Adapter configurations evaluated.} Two configurations were trained and compared:

\begin{itemize}
    \item \textbf{Standard adapter} (rank $r = 8$, \texttt{num\_layers = 16}): covers the upper 16 transformer layers --- 57\% of the 3B model and 100\% of the 1B model.
    \item \textbf{Fast adapter} (rank $r = 4$, \texttt{num\_layers = 8}): covers the upper 8 layers --- 29\% of the 3B model and 50\% of the 1B model. Prioritizes efficiency: fewer trainable parameters and faster inference.
\end{itemize}

The two configurations differ in both rank and layer count simultaneously, making it impossible to attribute performance differences to either factor in isolation; this is a known limitation of the comparison design.

\subsection{Synthetic Data Generation}

Following the LIMA hypothesis \cite{zhou2023lima} --- that a small number of carefully curated examples suffices for alignment --- we construct synthetic training corpora rather than collecting naturalistic patient data, motivated both by the impracticality of obtaining real pediatric clinical dialogues at scale and by ethical constraints. Training data was generated by Claude Opus 4.6 (Anthropic; March 2026) with Victoria Hospital's publicly available patient and family reference materials attached as grounding artifacts. Data is organized into five topic categories:

\begin{itemize}
    \item \textbf{What to Expect} --- medical procedures, equipment, and sensory experiences
    \item \textbf{Who Are These People?} --- age-appropriate explanations of clinical staff roles
    \item \textbf{Hospital Rules \& Routines} --- visiting hours, daily schedules, and preparation
    \item \textbf{Emotional Reassurance} --- responses to fear or confusion (non-diagnostic)
    \item \textbf{FAQs / General Curiosity} --- open-ended questions a child or caregiver might ask
\end{itemize}

Each adapter serves a single age group, with approximately 475 training examples per adapter across five categories. Safety constraints were enforced at generation time: the model was instructed never to imply a diagnosis, never to recommend medications, and to redirect emergency questions to clinical staff. A held-out validation set of 50 examples per age group was generated in a separate pass. Because both sets share the same generative process, perplexity measures in-distribution coherence rather than out-of-distribution generalization.

\subsection{Training Protocol}

Adapters are trained using the Adam optimizer with a cosine decay learning rate schedule: peak learning rate $1 \times 10^{-5}$, linear warmup over the first 100 steps, decaying to zero over 600 total steps. Batch size is 4; maximum sequence length is 768 tokens; LoRA dropout is 0.0. Loss is computed only over assistant response tokens (\texttt{mask\_prompt: true}). Both configurations are trained across both model sizes and both age groups, yielding 8 total runs (2 configurations $\times$ 2 model sizes $\times$ 2 age groups), all using seed 42.

\subsection{Evaluation Framework}

Three evaluation metrics are reported. \textbf{Readability:} five formula-based metrics (Flesch-Kincaid grade level \cite{kincaid1975fk}, SMOG \cite{mclaughlin1969smog}, Gunning Fog \cite{gunning1952fog}, Coleman-Liau \cite{coleman1975cli}, and type-token ratio) are computed over model-generated responses using the \texttt{textstat} library \cite{bansal2023textstat}. A pass/fail target of FK $\leq$ 7.0 is applied to the 5--11 adapter. \textbf{Latency and throughput:} per-response latency and tokens-per-second are measured sequentially on the Mac Mini M4; a 1.0-second average latency threshold reflects real-time VR interaction requirements. \textbf{Inter-role style separation:} a TF-IDF $+$ logistic regression classifier trained on the combined outputs of both adapters using 5-fold cross-validation (scikit-learn \cite{pedregosa2011sklearn}) quantifies whether the adapters produce lexically distinct registers; accuracy $\geq$ 0.90 is interpreted as strong separation.

\section{Results}

This section reports results for the eight training runs described in Section 3.5: two adapter configurations (Standard, $r = 8$, \texttt{num\_layers = 16}; Fast, $r = 4$, \texttt{num\_layers = 8}) across two model sizes (1B, 3B) and two age groups (5--11, 12--18). All results correspond to the final checkpoint at step 600. Base model (no-adapter) results are included as a zero-shot reference.

\subsection{Training Dynamics and Perplexity}

All four Fast adapter runs converge without divergence. Table~\ref{tab:res1} summarizes validation loss and derived perplexity ($e^{L_{\text{val}}}$) at the final step. Fast-3B achieves lower perplexity than Fast-1B for both age groups (5--11: 5.63 vs.\ 5.72; 12--18: 6.30 vs.\ 6.74). The 12--18 adapters show consistently higher perplexity than their 5--11 counterparts, consistent with more syntactically complex responses being harder to predict. Standard adapter perplexity requires a separate evaluation run and is not reported here.

\begin{table}[t]
  \caption{Fast adapter validation loss and perplexity at step 600.}
  \label{tab:res1}
  \small
  \begin{tabular}{llccc}
    \hline
    Config. & Age & Loss (step 1) & Loss (step 600) & PPL \\
    \hline
    Fast-1B & 5--11  & 3.211 & 1.744 & 5.72 \\
    Fast-1B & 12--18 & 3.245 & 1.908 & 6.74 \\
    Fast-3B & 5--11  & 3.023 & 1.728 & 5.63 \\
    Fast-3B & 12--18 & 2.998 & 1.840 & 6.30 \\
    \hline
  \end{tabular}
\end{table}

\subsection{Readability}

Table~\ref{tab:res2} reports readability metrics for all six configurations. All four fine-tuned adapters substantially outperform the base models on the FK $\leq 7.0$ target for the 5--11 group (72--84\% pass rate vs.\ 12--14\% for base models), driven largely by shorter response length (72--99 avg words vs.\ 152--210 for base models). The 12--18 adapters produce consistently higher FK grade levels than the 5--11 adapters (gap: 2.4--3.2 FK points), confirming that the two adapters have learned distinct communication registers.

A \textbf{configuration-ordering crossover} is visible in Table~\ref{tab:res2}: for the Standard configuration, Standard-3B outperforms Standard-1B on FK $\leq 7.0$ pass rate (84\% vs.\ 72\%); for the Fast configuration, this ordering reverses --- Fast-1B outperforms Fast-3B (82\% vs.\ 76\%). This crossover is consistent across all five readability metrics for the 5--11 group and is the central empirical observation of this evaluation.

\begin{table*}[t]
  \caption{Readability metrics (outputs from age-matched held-out prompts, 50 examples per group).}
  \label{tab:res2}
  \small
  \begin{tabular}{llccccccc}
    \hline
    Configuration & Age Group & FK Grade & SMOG & Gunning Fog & Coleman-Liau & TTR & Avg Words & FK $\leq$ 7.0 \\
    \hline
    Standard-1B & 5--11  & 6.1  & 8.8  & 7.7  & 6.7  & 0.712 & 74  & 36/50 (72\%) \\
    Standard-1B & 12--18 & 8.8  & 11.7 & 11.3 & 10.3 & 0.701 & 90  & ---          \\
    Standard-3B & 5--11  & 5.5  & 8.5  & 7.4  & 6.3  & 0.724 & 75  & 42/50 (84\%) \\
    Standard-3B & 12--18 & 8.7  & 11.4 & 10.9 & 10.0 & 0.714 & 92  & ---          \\
    Fast-1B     & 5--11  & 5.5  & 8.4  & 7.2  & 6.2  & 0.701 & 72  & 41/50 (82\%) \\
    Fast-1B     & 12--18 & 8.4  & 11.3 & 10.6 & 9.4  & 0.650 & 92  & ---          \\
    Fast-3B     & 5--11  & 5.9  & 8.9  & 7.9  & 6.6  & 0.636 & 99  & 38/50 (76\%) \\
    Fast-3B     & 12--18 & 8.3  & 11.2 & 10.7 & 9.7  & 0.613 & 122 & ---          \\
    Base-1B     & 5--11  & 9.5  & 12.1 & 12.2 & 9.0  & 0.508 & 201 & 7/50 (14\%)  \\
    Base-1B     & 12--18 & 10.8 & 13.1 & 13.4 & 11.2 & 0.527 & 210 & ---          \\
    Base-3B     & 5--11  & 9.4  & 12.1 & 12.2 & 9.6  & 0.609 & 152 & 6/50 (12\%)  \\
    Base-3B     & 12--18 & 10.6 & 12.9 & 13.3 & 11.4 & 0.565 & 206 & ---          \\
    \hline
  \end{tabular}
\end{table*}

\subsection{Latency and Throughput}

Table~\ref{tab:res3} reports inference latency for all configurations (sequential, no batching, Mac Mini M4). \textbf{Fast-1B is the only configuration meeting the 1.0-second average latency target} for the 5--11 adapter (0.93s avg, 70\% of responses under target). All 3B configurations exceed 2.3 seconds on average. Counterintuitively, Fast-3B is slower than Standard-3B despite fewer adapter parameters: Fast-3B generates $\sim$30\% more tokens per response (120/153 vs.\ 92/116 avg tokens), overriding any throughput benefit from the smaller adapter.

\begin{table*}[t]
  \caption{Inference latency and throughput (50 prompts per age group, sequential, Mac Mini M4).}
  \label{tab:res3}
  \small
  \begin{tabular}{llcccccc}
    \hline
    Config. & Age & Avg Lat. (s) & Min (s) & Max (s) & Avg Tok & Tok/s & $<$1.0s \\
    \hline
    Standard-1B & 5--11  & 1.09 & 0.66 & 1.69 & 89  & 81.7  & 21/50 (42\%) \\
    Standard-1B & 12--18 & 1.36 & 0.62 & 2.02 & 114 & 83.7  & 3/50 (6\%)   \\
    Standard-3B & 5--11  & 2.37 & 1.31 & 3.45 & 92  & 38.8  & 0/50 (0\%)   \\
    Standard-3B & 12--18 & 2.94 & 1.75 & 5.43 & 116 & 39.3  & 0/50 (0\%)   \\
    Fast-1B     & 5--11  & 0.93 & 0.50 & 1.40 & 88  & 94.3  & 35/50 (70\%) \\
    Fast-1B     & 12--18 & 1.20 & 0.75 & 1.79 & 115 & 95.4  & 12/50 (24\%) \\
    Fast-3B     & 5--11  & 2.83 & 1.31 & 6.80 & 120 & 41.9  & 0/50 (0\%)   \\
    Fast-3B     & 12--18 & 3.63 & 0.99 & 7.33 & 153 & 41.7  & 1/50 (2\%)   \\
    Base-1B     & 5--11  & 2.01 & 0.57 & 2.48 & 247 & 122.3 & 2/50 (4\%)   \\
    Base-1B     & 12--18 & 2.14 & 0.25 & 2.46 & 265 & 121.8 & 4/50 (8\%)   \\
    Base-3B     & 5--11  & 3.99 & 0.66 & 6.38 & 190 & 46.5  & 3/50 (6\%)   \\
    Base-3B     & 12--18 & 5.39 & 0.67 & 6.29 & 264 & 48.7  & 1/50 (2\%)   \\
    \hline
  \end{tabular}
\end{table*}

\subsection{Inter-Role Style Separation}

Table~\ref{tab:res4} reports TF-IDF + logistic regression classification accuracy (5-fold CV; chance = 0.50) for distinguishing 5--11 vs.\ 12--18 adapter outputs. All fine-tuned configurations achieve strong separation (0.89--0.94); base models reach only moderate separation (0.66--0.70), confirming that fine-tuning is the primary driver of distinct age registers. The same crossover observed in readability is present here: Standard-3B outperforms Standard-1B (0.920 vs.\ 0.890), while Fast-1B outperforms Fast-3B (0.940 vs.\ 0.900).

\begin{table}[t]
  \caption{Inter-role TF-IDF + LR classification accuracy (5-fold CV; chance = 0.50).}
  \label{tab:res4}
  \begin{tabular}{lc}
    \hline
    Configuration & Classification Accuracy \\
    \hline
    Standard-1B & $0.890 \pm 0.102$ \\
    Standard-3B & $0.920 \pm 0.051$ \\
    Fast-1B     & $0.940 \pm 0.073$ \\
    Fast-3B     & $0.900 \pm 0.077$ \\
    Base-1B     & $0.660 \pm 0.097$ \\
    Base-3B     & $0.700 \pm 0.063$ \\
    \hline
  \end{tabular}
\end{table}

\subsection{Discussion}

\textbf{Fast-1B is the recommended configuration} for VR deployment: it is the only configuration meeting the 1.0-second latency target (0.93s avg) while maintaining strong readability (FK 5.5 avg) and style separation (0.940 accuracy). The similarity between Standard and Fast on most quality metrics suggests that register adaptation for this task is achievable at rank 4 with 8 adapted layers, consistent with the LIMA-hypothesis framing: high-quality constrained training data reduces the adaptation capacity required to produce measurable style change.

The \textbf{crossover observation} --- Standard adapter better on 3B, Fast adapter better on 1B for both readability and style separation --- is the central empirical finding. We tentatively interpret this as a capacity-regularization effect: the 1B model's limited representational capacity may benefit from the implicit regularization of a smaller adapter, while the 3B model can exploit the additional degrees of freedom of rank 8. However, since rank and layer count co-vary across configurations and all runs use a single seed, this crossover is not statistically confirmed and requires a controlled rank sweep with multiple seeds to establish.

Three structural limitations apply: (1) single-seed runs --- no variance estimates for any metric; (2) no Standard adapter perplexity --- training logs were not retained; (3) in-distribution evaluation only --- both training and validation data share the same generative process.

\section{Conclusion}
Conclusion

\section*{Appendix}

\subsection{Repository Link}
% [TODO: add repository URL]
\noindent [To be added.]

\subsection{Group Member Contributions}
% [TODO: fill in]
\noindent [To be added.]

\subsection{Claude Transcripts}
% [TODO: specify Claude version and session details]
\noindent [To be added.]

\clearpage
\begin{thebibliography}{99}

\bibitem{zhou2023lima}
C.~Zhou, P.~Liu, P.~Xu, S.~Iyer, J.~Sun, R.~Maddela, X.~Peng, S.~Shrivastava, M.~Lewis, L.~Zettlemoyer, and O.~Levy, ``LIMA: Less Is More for Alignment,'' in \textit{Advances in Neural Information Processing Systems (NeurIPS)}, 2023.

\bibitem{dubey2024llama3}
A.~Dubey et al., ``The Llama 3 Herd of Models,'' \textit{arXiv preprint arXiv:2407.21783}, Meta AI, 2024.

\bibitem{apple2023mlx}
Apple, ``MLX: Efficient and Flexible Machine Learning on Apple Silicon,'' 2023. [Online]. Available: https://github.com/ml-explore/mlx

\bibitem{hu2022lora}
E.~J.~Hu, Y.~Shen, P.~Wallis, Z.~Allen-Zhu, Y.~Li, S.~Wang, L.~Wang, and W.~Chen, ``LoRA: Low-Rank Adaptation of Large Language Models,'' in \textit{Proc. Int. Conf. Learning Representations (ICLR)}, 2022.

\bibitem{kincaid1975fk}
J.~P.~Kincaid, R.~P.~Fishburne, R.~L.~Rogers, and B.~S.~Chissom, ``Derivation of New Readability Formulas for Navy Enlisted Personnel,'' Research Branch Report 8-75, Naval Technical Training Command, 1975.

\bibitem{mclaughlin1969smog}
G.~H.~McLaughlin, ``SMOG Grading --- A New Readability Formula,'' \textit{J. Reading}, vol.~12, no.~8, pp.~639--646, 1969.

\bibitem{gunning1952fog}
R.~Gunning, \textit{The Technique of Clear Writing}. New York, NY, USA: McGraw-Hill, 1952.

\bibitem{coleman1975cli}
M.~Coleman and T.~L.~Liau, ``A Computer Readability Formula Designed for Machine Scoring,'' \textit{J. Applied Psychology}, vol.~60, no.~2, pp.~283--284, 1975.

\bibitem{bansal2023textstat}
S.~Bansal, ``textstat: Python Library to Calculate Readability Statistics of a Text Object'' (v0.7.13), 2023. [Online]. Available: https://pypi.org/project/textstat/

\bibitem{pedregosa2011sklearn}
F.~Pedregosa et al., ``Scikit-learn: Machine Learning in Python,'' \textit{J. Machine Learning Research}, vol.~12, pp.~2825--2830, 2011.

\end{thebibliography}

\end{document}
